\section{Memory, Cognition, and Behavioral Continuity}
\label{sec:memory-cognition}

GAWorld separates persistence from generation.  An agent's continuity is not
held solely in an LLM context window: structured state, episodic records,
habits, intentions, relationships, environmental preferences, growth profiles,
and skills are written to mechanism-specific stores and reintroduced at later
steps.  Figure~\ref{fig:memory-growth} summarizes this path.  The resulting
system can preserve and expose information across simulated days; whether its
memory dynamics resemble those of people is a separate empirical question.

\begin{figure}[t]
  \centering
  \includegraphics[width=0.94\linewidth]{memory_growth.pdf}
  \caption{Persistent experience and growth paths.  Tick-level episodes feed
  retrieval and rule-based updates; configured day-boundary processes may
  consolidate episodes into semantic memories or private skills.  Solid
  storage paths denote implemented persistence, not validated models of human
  remembering or learning.}
  \label{fig:memory-growth}
\end{figure}

\subsection{Episodic persistence and retrieval}

The basic store retains an agent's memory list, schedule, available actions,
locations, and location--action biases as JSON, alongside append-only textual
logs and simulation state.  A second experience layer records episodes in
JSONL and persists habits, daily intentions, relationships, and environmental
preferences in separate per-agent files.  This division supports inspection
and selective restart, although a complete restart requires all of these
stores rather than a single canonical agent file.

For content-addressed recall, memory entries are also indexed in SQLite with
type, simulated time, salience, recall count, and an embedding.  The configured
embedding path may call an embedding provider; if that path is unavailable,
the implementation falls back to a deterministic hash embedding.  Retrieval
combines vector similarity with salience, age decay, and prior recall, then
optionally boosts entries related to the agent's current growth interests.
Thus ranking is a specified scoring procedure over stored entries, whereas the
text being ranked may have originated in an LLM-generated action, reflection,
diary, news item, or external source.

The day-boundary lifecycle adds two optional transformations.  Consolidation
selects recent episodic rows and asks an LLM to distill them into a small set of
semantic entries; malformed output or insufficient material produces no new
entry.  Decay is LLM-free: after a configured period it lowers the salience of
old, unrecalled episodic rows and deletes rows already below a floor, while
protecting semantic, external-information, and procedural entries.  A separate
episode store also computes exponential salience decay and enforces a size cap.
These two decay paths operate over different stores and should not be mistaken
for one unified cognitive theory.

External retrieval-augmented memory is likewise optional.  Agent-derived
queries can retrieve web material, after which an LLM may assess or compress
it into external-information entries.  This path depends on a caller-supplied
fetch adapter and degrades to no operation on failure.  Consequently, an
experimental report must disclose whether external retrieval, provider-based
embeddings, consolidation, and decay were enabled.

\subsection{Intentions, habits, and reactive behavior}

Behavioral continuity combines model-mediated interpretation with explicit
state-transition rules.  At day boundaries, recent episodes and current state
can be supplied to an LLM to produce priorities, avoidances, social targets,
and recovery goals.  When the call budget is exhausted or parsing fails, a
rule-based fallback constructs intentions from state and recent experience.
Habit strength, relationship signals, physiological and time-pressure needs,
episode salience, and many mood or stress adjustments are updated by bounded
numeric rules.  These values enter subsequent prompt context rather than being
left as unstructured prose.

The dynamic-behavior module evaluates candidate interruptions from needs,
co-location, incoming messages, environmental events, local physical
conditions, cascades, and spontaneous urges.  Acceptance probabilities depend
on event intensity, activity commitment, personality proxies, relationship
closeness, and agent state.  Candidate generation and schedule insertion are
LLM-free but stochastic; they are repeatable only when the relevant random
state is seeded.  Once an interruption wins, the ordinary activity-adjustment
or action path may invoke an LLM to phrase or refine the response.  GAWorld
therefore distinguishes rule-based opportunity and constraint generation from
model-generated interpretation and action text.

Spatial preference illustrates continuity without language generation.  A
crowd surge or other location anomaly raises a time-bucketed aversion score;
scores decay with recency, and a later destination can be redirected when its
aversion crosses a configured threshold.  This mechanism implements learned
avoidance as a transparent rule.  It does not establish that the learning rate,
threshold, or decay reproduces human spatial choice.

\subsection{Interests, skills, and productive activity}

Growth profiles contain bounded hobby and skill items with phase, engagement,
competence, practice history, and activity templates.  Their initial content
may be derived by an LLM and cached by a profile signature; deterministic
fallback profiles are available when generation fails.  Episode matching,
practice updates, phase classification, daily retention decay, retirement of
stale interests, and socially mediated adoption then use explicit rules and a
seedable random generator.  The labels borrow a developmental vocabulary, but
their implementation is a simulation mechanism rather than a validated model
of human interest development.

The skill subsystem maintains a curated global library and per-agent private
skills.  Available skills can be injected into cognition context, while an
optional LLM consolidation step extracts a reusable procedure from recent
episodes and saves it as a private Markdown skill.  Parsing, schema checks, and
registry precedence constrain this output; they do not demonstrate that the
procedure is correct or transferable outside the generated setting.

Productive activity extends this pattern to a mock work market.  Cached,
LLM-derived capability profiles and current growth items inform routing.  A
JSONL queue provides claim and recovery semantics, a worker pool dispatches
briefs concurrently, and code, content, teaching, and web-design adapters use
an LLM to produce files.  Lightweight validators check properties such as
Python syntax, Markdown structure, or HTML/SVG parseability before successful
results are ingested into agent state and market settlement.  These checks and
persisted artifact paths demonstrate a functioning production pipeline; they
are not evidence that an artifact meets professional, pedagogical, factual, or
security standards.
